\documentclass[11pt]{article}
\usepackage[a4paper, margin=0.4in]{geometry}
\usepackage{trimclip}
\usepackage{tabularx}
\usepackage[outline]{contour}% http://ctan.org/pkg/contour
\usepackage{tabto}
\usepackage{blindtext}
\usepackage{hyperref} \usepackage[dvipsnames]{xcolor}
\usepackage{multicol}
\usepackage{etoolbox}
\usepackage{accsupp}
\usepackage[fixed]{fontawesome5}
\usepackage{ifxetex,ifluatex}
\usepackage{scrlfile}
\usepackage{xparse}
\setlength{\columnseprule}{1pt}
\usepackage{helvet}

% \usepackage{sansmathfonts}
% \usepackage[T1]{fontenc}

\usepackage{DejaVuSansCondensed}
\renewcommand*\familydefault{\sfdefault} %% Only if the base font of the document is to be sans serif
\usepackage[T1]{fontenc}

\definecolor{Aquamarine}{rgb}{0.11, 0.34,0.45 }
\def\columnseprulecolor{\color{Aquamarine}}
\newcommand{\itemmarker}{{\small\textbullet}}
\newcommand{\ratingmarker}{\scriptsize\faCircle}
\setlength{\parindent}{0cm}
\renewcommand{\baselinestretch}{1.06}

\hypersetup{
    colorlinks=true,
    linkcolor=blue,
  filecolor=magenta,
  urlcolor=Aquamarine,
  pdftitle={cv},
    pdfpagemode=FullScreen,
    }

\newenvironment{info}[2]
{
  \textbf{#1}: \tabto{1.5in}{#2}
}{
  \smallbreak
}

\newenvironment{til}[2]
{
  \large
  \textcolor{Aquamarine}{\textbf{#1}}
  \vspace{0.05in}
  \textcolor{Aquamarine}{\hrule}
}{
  \vspace{0.05in}
}

\newenvironment{itemexp}[5]{
  \contourlength{0.22pt}
  \contournumber{10}%
  \contour{black}{#1} \hfill \begin{scriptsize}#2\end{scriptsize}\\
  \scriptsize\textit{#3} \hfill \textit{#4}\\\normalsize
  \normalsize #5}{\smallbreak}

\newcommand{\cvskill}[2]{%
  \textcolor{black}{\small{#1}}\hfill
  \BeginAccSupp{method=plain,ActualText={#2}}
  \foreach \x in {1,...,5}{%
    \ifdimequal{\x pt - #2 pt}{0.5pt}%
    {\clipbox*{0pt -0.25ex {.5\width} {\totalheight}}{\color{Aquamarine}\ratingmarker}%
     \clipbox*{{.5\width} -0.25ex {\width} {\totalheight}}{\color{Aquamarine!30}\ratingmarker}}
    {\ifdimgreater{\x bp}{#2 bp}{\color{Aquamarine!20}}{\color{Aquamarine}}\ratingmarker}%
  }\EndAccSupp{}
  \smallbreak\par%
}

\newcommand{\cvcertificate}[2]{%
  \small{\textcolor{black}{\footnotesize{#1}}}\smallbreak
  \tiny {#2}
  \smallbreak\par%
}

\newcommand{\subtitle}[2]{
  \textcolor{Aquamarine}{\textbf{#1}} \textcolor{Aquamarine}{\faIcon[solid]{#2}}\smallbreak
}{}

\newcommand{\subsubtitle}[1]{
  \textcolor{Aquamarine}{\small\textbf{#1}}\smallbreak
}{}

\pagestyle{empty} % Suppress page numbers
\usepackage{eso-pic}
\usepackage{tikz}
\usepackage{fancyhdr}
\usepackage[explicit,]{titlesec}

\begin{document}
\medbreak
\huge
\contourlength{0.22pt}
\contournumber{10}%
\contour{black}{Roberto Andrés Alvarado Moreira}
\medbreak
\footnotesize
\normalsize
Desarrollador Senior de I\&D enfocado en verificación de chips, C++20 e innovación aplicada en IA.

\vspace{-0.2cm}
\small
\textcolor{Aquamarine}{
  \begin{center}
    \begin{tabularx}{0.8\textwidth} { 
       >{\raggedright\arraybackslash}X 
       >{\raggedright\arraybackslash}X}
        \faIcon{envelope}{ robdres123@gmail.com} &\faIcon{github}
        {\normalsize\url{https://github.com/Robdres}}\\[0.1cm]
        \faIcon{mobile}{ (+593)988718378} &\faIcon{map} {\small Quito, Ecuador}\\
    \end{tabularx}
  \end{center}
}
\normalsize
\medbreak
\vspace{0.05cm}
\normalsize
\begin{minipage}[t]{0.3\textwidth}
  \begin{til}{Herramientas}{1}
  \end{til}
  \medbreak
  \footnotesize
  \subtitle{Lenguajes}{globe}
  \medbreak
  \cvskill{Español}{5}
  \cvskill{Inglés}{5}
  \medbreak
  \subtitle{Lenguajes de Programación}{laptop}
  \medbreak
  \cvskill{C++20}{5}
  \cvskill{Python}{5}
  \cvskill{Rust}{3}
  \cvskill{C}{3}
  \cvskill{Bash}{3}
  \cvskill{MySQL}{3}
  \cvskill{Latex}{3}
  \medbreak
  \subtitle{IA}{robot}
  \medbreak
  \cvskill{Copilot}{4}
  \medbreak
  \subtitle{Herramientas de desarrollo}{keyboard}
  \medbreak
  \cvskill{Vim}{5}
  \cvskill{Figma}{4}
  \cvskill{VSCode}{3}
  \medbreak
  \subtitle{Plataformas}{ubuntu}
  \medbreak
  \cvskill{Linux}{5}
  \cvskill{Docker}{3}
  \cvskill{Windows}{3}
  \medbreak
  \subtitle{Frameworks y Librerías}{database}
  \medbreak
  \cvskill{Angular}{5}
  \cvskill{TensorFlow}{4.5}
  \cvskill{NodeJs}{4}
  \medbreak

  \begin{til}{Habilidades}{1}
  \end{til}
  \medbreak
  \cvskill{Liderazgo}{5}
  \cvskill{Pensamiento crítico}{5}
  \cvskill{Trabajo en Equipo}{4}
  \cvskill{Comunicación}{3.5}
\vspace*{\fill}
\end{minipage}
\textcolor{NavyBlue}{\hfill\vline \hfill}
\begin{minipage}[t]{0.61\textwidth}
  \begin{til}{Experiencia Laboral}{1}
  \end{til}
  \smallbreak
  \begin{itemexp}
    {Desarrollador Senior de I\&D}{Junio,2023-Presente}{Synopsys}{Santiago, Chile}
    {Desarrollé y mantuve algoritmos de verificación de chips. Entregué 5 proyectos de I\&D de escala completa y lideré la innovación de IA en Formality, logrando 60\% de uso activo y 90\% de adopción en pruebas entre desarrolladores.}
  \end{itemexp}
  \smallbreak
  \begin{itemexp}
    {Práctica Quant}{Mayo,2023-Junio,2023}{Ney Torres Investment Fund}{Quito, Ecuador}
    {Analicé datos financieros y construí una herramienta comparativa para métricas de trading.}
  \end{itemexp}
  \smallbreak
  \begin{itemexp}
    {Ingeniero Fullstack Junior}{Enero,2022-Febrero,2023}{Tecsicom}{Quito, Ecuador}
    {Construí y entregué una plataforma basada en Angular para la operación diaria de una importante empresa de marketing en Ecuador.}
  \end{itemexp}
  \smallbreak
  \begin{itemexp}
    {Asistente de cátedra}{Junio,2022-Marzo,2023}{Universidad
    San Francisco de Quito}{Quito, Ecuador}
    {Apoyé actividades académicas junto al profesor David Hervas, PhD en Matemáticas.}
  \end{itemexp}
  \smallbreak
  \begin{til}{Educación}{1}
  \end{til}
  \smallbreak
  \begin{itemexp}
    {Maestría en Inteligencia Artificial Aplicada}{Abril,2025-Abril,2026}{Universidad
    Tecnica Particular de Loja}{Loja, Ecuador}
    {\textit{Cursos}: Redes Neuronales, Sistemas de Clasificación, Industria 4.0}
  \end{itemexp}
  \smallbreak
  \begin{itemexp}
    {Licenciatura en Matemáticas}{Enero,2019-Mayo,2023}{Universidad
    San Francisco de Quito}{Quito, Ecuador}
    {\textit{Cursos}: Análisis Numérico, Redes Neuronales, Análisis Real}
  \end{itemexp}
  \smallbreak
  \begin{itemexp}
    {Minor en Filosofía}{Enero,2018-Mayo,2024}{Universidad
    San Francisco de Quito}{Quito, Ecuador}
    {\textit{Cursos}: Filosofía de la religión, Historia de la filosofía, Filosofía del Proceso,
    Diálogos Platónicos}
  \end{itemexp}
  \smallbreak
  \begin{til}{Proyectos}{1}
  \end{til}
  \smallbreak
  \begin{itemexp}
    {Testores}{Agosto,2021-Mayo,2023}{Universidad
    San Francisco de Quito}{Quito, Ecuador}
    {Trabajé con métodos de testores (rough sets) para reducir atributos de clasificación y mejorar tiempos de procesamiento.}
  \end{itemexp}
  \begin{itemexp}
    {Palm Tree}{Enero,2023-Mayo,2023}{}{Quito, Ecuador}
    {Creé métricas de desempeño para dispositivos infrarrojos usando datos de palma en colaboración con NC State.}
  \end{itemexp}
\end{minipage}
\end{document}
